% Seccion: Principios de transferencia de calor
\section{Principios de transferencia de calor}
\label{sec:cap02_principios_transferencia}

La conducción de calor en medios sólidos está regida por la ley de Fourier, donde el flujo es proporcional al gradiente de temperatura.
En el colector, la conducción ocurre principalmente a través de la chapa absorbedora y los espesores del aislante de poliuretano expandido.
La convección de calor entre la placa y la corriente de aire está descrita por la ley de enfriamiento de Newton.
El coeficiente convectivo depende fuertemente de las propiedades termo-físicas del aire y de la velocidad del flujo.
La transferencia por radiación de onda larga entre la placa absorbedora caliente y la cubierta de vidrio es un mecanismo importante de pérdida.
Esta radiación sigue la ley de Stefan-Boltzmann corregida por la emisividad e inductancias de las superficies involucradas.

\begin{cajainfo}{Ecuación Convectiva de Newton}
El flujo de calor convectivo útil se expresa matemáticamente mediante la siguiente ecuación:
q_c = h \cdot (T_p - T_f)
Donde h es el coeficiente de transferencia de calor, Tp es la temperatura de la placa y Tf la del fluido.
\end{cajainfo}
